\section{Introduction}

Deep learning is now the standard approach for segmenting anatomical structures
in medical images, and its outputs increasingly guide clinical decisions --
from surgical planning to disease
monitoring~\cite{isensee2021nnunet,isensee2024revisited}. In cervical-spine MRI,
automated segmentation of vertebral bodies and intervertebral discs supports
stenosis grading and degeneration tracking~\cite{zhou2025cspineseg}. If such
models work worse for certain patient groups -- differing by sex, race, or age --
they risk worsening the health disparities already present in spine care.
With the EU AI Act classifying medical AI as high-risk and requiring
documented demographic representativeness and bias
assessment~\cite{euaiact2024}, fairness audits are no longer optional: they are
a precondition for deployment.

Fairness is well studied for classification
tasks~\cite{bird2020fairlearn,eeoc1978uniform}, but medical image segmentation
raises challenges with no classification analogue. First, the output is a spatial
map of millions of voxels, not a single decision, so there is no standard
fairness metric: one must reduce the map to a single number (e.g.\ Dice or HD95),
and that reduction can itself introduce bias when anatomy differs systematically
between demographic groups~\cite{raina2023ndsc,tian2024fairseg}. Second,
segmentation annotations are far more expensive and subjective than classification
labels, with inter-annotator disagreement that can vary systematically across
patient groups~\cite{parikh2025labelbias}. The high cost of expert annotation forces
datasets to mix expert (gold) and machine-generated (silver) labels -- and if the
silver labels carry demographic biases of their own, the reference used to judge
fairness becomes unreliable. Parikh et al.~\cite{parikh2025labelbias} call this
the biased ruler effect.

Despite growing interest in fairness for medical imaging, segmentation remains
underexplored. Tian et al.~\cite{tian2024fairseg} released Harvard-FairSeg -- the
first fairness dataset for medical segmentation (ophthalmology, ICLR 2024) --
showing how recently the field acquired basic infrastructure. In breast-MRI
segmentation, Parikh et al.~\cite{parikh2025labelbias,parikh2025whofails} showed
that machine-generated reference labels can inflate observed disparities and that
training on biased silver labels widens demographic gaps. The FairMedFM
benchmark~\cite{jin2024fairmedfm} tested foundation-model fairness across 17
datasets and found widespread bias, but its only spine entry (SPIDER, lumbar) was
limited to sex and did not correct for volume differences. Puyol-Ant\'{o}n et
al.~\cite{puyolanton2021fairness} and Danaee et
al.~\cite{danaee2026brain} audited cardiac and brain segmentation,
respectively, but neither covered spine. Gichoya et
al.~\cite{gichoya2022reading} showed that AI models can infer patient race
directly from medical images, and Petersen et
al.~\cite{petersen2023feature} showed that encoding a protected attribute does not
necessarily cause a performance gap -- a distinction we revisit here. To the best of our knowledge,
no dedicated fairness audit exists for cervical-spine MRI segmentation
across multiple demographic axes, and no prior work has tested whether label
provenance -- in training or evaluation -- changes fairness verdicts in this
anatomy.

\smallskip\noindent\textbf{Contributions.}\enspace
We present the first demographic fairness audit of cervical-spine MRI
segmentation, evaluating performance across sex, race, and age on the CSpineSeg
dataset~\cite{zhou2025cspineseg} ($1{,}254$ exams). Our main findings are:
\begin{enumerate}
  \item The model is fair across all three protected attributes by every
    standard metric and threshold tested.
  \item Replacing the expert reference labels with machine-generated ones flips
    the fairness verdict for age -- revealing a ``false confidence'' mode
    of biased-ruler bias that complements the ``false magnitude'' mode of Parikh
    et al.~\cite{parikh2025labelbias}.
  \item Training on machine-generated labels does not amplify demographic bias in
    this anatomy, unlike the amplification observed in breast
    MRI~\cite{parikh2025labelbias}.
  \item Encoder probing confirms that the observed age gradient reflects
    intrinsic anatomical difficulty rather than algorithmic harm.
\end{enumerate}
